  \intro


Aqu\'{\i} la introducci\'on del trabajo. En \cite{guidedissertation} pueden leerse
estos consejos:

The dissertation should be divided into chapters and sections appropriate to
the topic and type of dissertation chosen. The following elements are typical
of the traditional dissertation. You should discuss the overall structure of
your dissertation with your academic supervisor.

The Introduction to the dissertation should set out the background to the
research study and address the following areas:

\medskip
\noindent\textsl{The context in which the research took place}
\begin{itemize}
 \item What is the background, the context, in which the research took place?
 \item Why is this subject or issue important?
 \item Who are the key participants and/or `actors' in the area under investigation?
 \item Are there important trends or pivotal variables of which the reader
       needs to be made aware?
 \item A clear and succinct statement of the aims and objectives that the
       dissertation is going to address.
 \item Have you presented a clear and unambiguous exposition of your research
       aim, the objectives you will address to meet this aim and your research
       questions?
\end{itemize}

\medskip
\noindent\textsl{The reasons why this study was carried out}
\begin{itemize}
 \item Was this study undertaken for example in order to test some aspect of
professional or business practice or theory or framework of analysis?
 \item Was the research carried out to fulfil the demands of a business
organisation?
\end{itemize}

\medskip
\noindent\textsl{The way the Dissertation is to be organised}

You should write your dissertation with the idea in mind that the intended reader and
reviewer has some shared understanding of the area being investigated, however,
underpinning concepts and arguments still need to be included as otherwise the depth
of research will be compromised. In this way, you will not be tempted to make too
many implicit assumptions, i.e. by making the erroneous assumptions that the reader
has your degree of knowledge about the matters in question or can follow, exactly,
your thought processes without your spelling them out. It should be a document which
is `self-contained' and does not need any additional explanation, or interpretation, or
reference to other documents in order that it may be fully understood.

This short final section of the Introduction should tell the reader what topics are going
to be discussed in each of the chapters and how the chapters are related to each other.
In this way, you are, in effect, providing the reader with a `road map' of the work
ahead. Thus, at a glance, they can see (1) where they are starting from, (2) the context
in which the journey is taking place, (3) where they are going to end up, and (4) the
route which they will take to reach their final destination. Such a `map' will enable
the reader to navigate their way through your work much more easily and appreciate
to the maximum what you have done.
